Processing factors (solar vs. freeze-drying; steam-blanching vs. no pretreatment) significantly determine both initial micronutrient levels and changes in mineral contents during storage of cocoyam leaves. However, the study showed that process factor did not influence nutrient stability alone, interactions (drying x pretreatment x temperature x time) were more important. In general, drying changed the matrix of cocoyam leaves from a water-rich, dilute one to a solid-dense matrix. As a consequence, the freshly dried samples had higher percentages of minerals than fresh leaves. When comparing solar and freeze drying it can be generally said that the two drying methods in combination with steam blanching were differently effective for the minerals. The samples that recorded the highest phosphorus content was the solar-dried, steam-blanched sample, whereas calcium contents were the highest in freeze-dried, unblanched samples.  
Similar effects can be observed after steam blanching that increased initial mineral content, especially of calcium and phosphorus. But also here it must be said that blanching interacts with drying and temperature, meaning its benefits depend on the full process chain.
